% Appendix B: Proofs and finite counts of Part II
\label{app:proofs-body}

This appendix supplies the details omitted from Part~II: the adequacy induction, the
countermodels that force each hypothesis of
Definition~\ref{def:adequacy-hypotheses}, the critical-pair configurations, the monotonicity
computation behind strong normalization, and the envelope incidence count.

\subsection{The adequacy induction in detail}

We record the rule-by-rule obligations of Theorem~\ref{thm:adequacy}.

\begin{center}
\begin{tabular}{@{}p{0.20\textwidth}p{0.30\textwidth}p{0.42\textwidth}@{}}
\toprule
Rule & Obligation & Discharged by \\
\midrule
assumption & $\rho$ realises the thread & Definition~\ref{def:realisation} \\
$\wedge$I & family intersection is inhabited & hypothesis (3) \\
$\wedge$E & projection of an inhabited family & hypothesis (3) \\
$\to$I & discharged thread replaced by witness & hypothesis (1) \\
$\to$E & consumed thread matched & hypotheses (1), (2) \\
$\forall$I & a witness slot may be chosen & hypothesis (1) \\
$\forall$E & instantiation stays inside the family & hypothesis (3) \\
atom & witness or emptiness certificate & hypothesis (4) \\
\bottomrule
\end{tabular}
\end{center}

The completeness direction uses hypothesis (4) twice: once to build the canonical assignment
and once to certify emptiness when the goal family is empty.  This is why (4) cannot be
weakened to "either a witness or nothing found": without the certificate an empty family
would be indistinguishable from an unfinished search, and by
Warning~\ref{warn:empty-unresolved} the two must not be conflated.

\subsection{Countermodels for the four hypotheses}

Each hypothesis is necessary for the equivalence of Theorem~\ref{thm:adequacy} in the
following exact sense: there is a fixture on which the equivalence fails, the other three
hypotheses hold, and the failure is attributable to the dropped hypothesis.

\begin{enumerate}
\item \textbf{Without (Cartesian).}  Take a carrier whose slot domains are dependent, so a
  partial assignment may admit no extension.  Then an implication whose premise is
  satisfiable in isolation but not jointly with the context is derivable by the rule while
  its conclusion has empty family.
\item \textbf{Without (Compatibility).}  Take three slots with two constraints such that each
  pair is compatible but the triple is not.  Then pairwise sections are all nonempty while
  the joint family is empty --- the exact situation of
  Proposition~\ref{prop:triadicity-test}, read as a failure of adequacy rather than as
  evidence for triadicity.
\item \textbf{Without (Transformer).}  Take a connective with a transformer that is sound but
  not complete; then a derivable conclusion has nonempty family but the transformer's image
  is strictly larger, so family inclusion does not track derivability.
\item \textbf{Without (Exhaustiveness).}  Take a carrier where some atomic family is empty
  with no certificate; the completeness direction fails at that atom, and the failure is
  reported as \emph{unresolved} rather than as falsity.
\end{enumerate}

\subsection{The four critical-pair configurations}

For two overlapping redexes (Definition~\ref{def:overlapping-redexes}) the ledger blocks
compare as follows.  Write the outer detour's ledger as $L\,[x^{b}]\,Q\,S$ and the inner one
as sitting either inside $Q$ or inside $S$.

\begin{center}
\begin{tabular}{@{}clp{0.46\textwidth}@{}}
\toprule
Case & Position of the inner redex & Relation of the two reduction orders \\
\midrule
(1) & in $Q$, on the discharged spine & joined after one splice step \\
(2) & in $Q$, off the spine & commute exactly \\
(3) & in $S$, adjacent block & joined after one splice step \\
(4) & in $S$, order-adjacent block & commute exactly \\
\bottomrule
\end{tabular}
\end{center}

In cases (1) and (3) the two-step results are related by
Theorem~\ref{thm:substitution-splice} applied to the nested substitution; the freshness
condition is inherited from the outer redex.  In cases (2) and (4) the blocks are disjoint, so
the transports of Theorem~\ref{thm:root-transport} act on independent subterms.

\subsection{Monotonicity of the ledger counts}

Recall that $\Lambda_D$ is the complete proof-tree ledger and that $\operatorname{dchg}$ counts
its discharged records (Definition~\ref{def:proof-ledger}).

\begin{enumerate}
\item A root transport retires exactly the bound target record $x$:
  $\Lambda_R=L\circ[x^{\mathrm{discharged},b}]\circ Q\circ\Lambda_\sigma$ becomes
  $\Lambda_C=L\circ\Lambda_\sigma\circ Q$, with no other record deleted, duplicated,
  relabelled, or renamed.  Both $|\Lambda|$ and $\operatorname{dchg}$ therefore fall by one,
  and the two subterms $Q$ and $S$ are unchanged as trees.
\item An introduction ancestor above the redex discharges the same surviving open record on
  both sides, so it raises both counts equally; either elimination ancestor concatenates the
  same sibling ledger on both sides.  A step strictly inside a context composes with these
  ancestors, and induction along the proof-position path gives the same two decrements at
  the outer root.
\item In the unfolded representation the node count also drops by exactly three, because
  linearity grafts the argument exactly once: the reduct removes the elimination node, the
  introduction node, and the target assumption leaf.
\end{enumerate}

Since $\N$ is well founded, an infinite sequence of certified steps would give an infinite
strictly descending sequence of natural numbers, which is
Theorem~\ref{thm:strong-normalization}.

Linearity enters at the first item and at the node count.  With contraction available the
argument fails, and it fails through the ledger and the graft rather than through a detour
count: see the remark of \S\ref{sec:normalization-body} showing that contracting one
principal detour can expose another.

\subsection{The ledger splice}

For a context $L\,Q\,R$ with $Q$ the transported middle block, the ledger of the composite is
$L\circ\Lambda_\sigma\circ R$ by Theorem~\ref{thm:substitution-splice}.  Concretely, if
$\Lambda_\sigma$ replaces the $i$-th antecedent occurrence by a block of length $\ell_i$, then
the composite ledger has the same left and right blocks and a middle block of length
$\sum_i\ell_i$; freshness is precisely the condition that no occurrence of the original block
appears twice in the image.

\subsection{Envelope incidence count}

The truncated octahedron has $24$ vertices, all of degree $3$, and faces consisting of six
squares and eight hexagons.  Then
\[
  \#\text{face--vertex incidences}=6\cdot4+8\cdot6=24+48=72=24\cdot3,
\]
which agrees with the vertex count.  The relation boundary count of
Proposition~\ref{prop:envelope} is $2m=6$ for $m=3$, matching
Theorem~\ref{thm:counting-bound}.
